\subsection{Scalar normalization at a specified output radius}
\label{sec:scalar-rms-amplitude}

We retain the output radius as a parameter. A larger radius reduces
the requested sign mean and permits an initial thinning gate.
The geometric-majority identity in Lemma~\ref{lem:rmsgeommajority}
gives the upper bound in the scalar law below.

\begin{lemma}[A second-derivative bound from uniformly finite mean work]
\label{lem:new-curvature-transfer}
Let $f\colon[-1,1]\to[-1,1]$ be twice continuously differentiable.
Suppose a single exact factory, valid for every source-sign mean
$z\in[-1,1]$, returns a sign of mean $f(z)$ and uses at most $K$
source requests in expectation uniformly over $z$.  Then
\[
 K\ge\frac12\sup_{-1<t<1}(1-t^2)|f''(t)|.
\]
Only the first moment of the source count is assumed.
\end{lemma}
\paragraph{Proof idea.}
Run the factory on uniformly sampled entries of a finite hidden
vector. The finite transcript curvature bound applies to this experiment
without a second moment for the stopping time. Second derivatives of its
Bernstein polynomial converge to the requested scalar derivative.

\begin{proof}
For an integer $m\ge2$, give a query algorithm a hidden sign vector
$x\in\{-1,1\}^m$.  Simulate each requested source sign by drawing an
independent uniform index in $[m]$ and returning its coordinate.  Repeated
coordinates can be cached.  Conditional on $x$, these source signs are
independent and have mean $M(x)=m^{-1}\sum_i x_i$.  Thus the output
has exact conditional mean $f(M(x))$, and its number $Q$ of distinct
input probes satisfies $\mathbb E[Q\mid x]\le K$.

Give the hidden coordinates independent signs of common mean
$t\in(-1,1)$ and put $H_m(t)=\mathbb E_t f(M)$. The finite-input
transcript bound of Lemma~\ref{lem:transcript} gives
\[
 |H_m''(t)|\le\frac{2\mathbb E_tQ}{1-t^2}
             \le\frac{2K}{1-t^2}.
\]
It requires no moment assumption on the stopping time beyond the
stated mean-work bound.

It remains to pass to the desired scalar function.  With
$p=(1+t)/2$, the function $H_m$ is the Bernstein polynomial
\[
 H_m(t)=\sum_{j=0}^m\binom mjp^j(1-p)^{m-j}f(2j/m-1).
\]
Let $J$ have distribution $\operatorname{Bin}(m-2,p)$, and let
$U,V$ be independent uniforms on $[0,1]$, independent of $J$.
Twice differentiating the polynomial, then expressing its second
finite difference by two integrals, gives
\[
 H_m''(t)=\left(1-\frac1m\right)
 \mathbb E\,f''\!\left(\frac{2J}{m}-1+\frac{2(U+V)}m\right).
\]
The argument lies in $[-1,1]$ and converges to $t$ in probability,
since its variance is $O(1/m)$ and its mean differs from $t$ by
$-2t/m$.  Uniform continuity and boundedness of $f''$ imply
$H_m''(t)\to f''(t)$.  Pass to the limit in the preceding inequality
and then take the supremum over $t$.
\end{proof}

\begin{theorem}[Scalar RMS cost at a specified amplitude]
\label{thm:scalar-rms-amplitude}
Let $\kappa>0$ and $A\ge1$.  Consider factories valid for every
unknown sign mean $z\in[-1,1]$ whose output mean is
\[
 f_{\kappa,A}(z)
 =\frac{z}{A\sqrt{\kappa^{-1}+z^2}}.
\]
Let $Q^*(\kappa,A)$ be the infimum, over all such exact factories,
of their worst-case expected number of source requests.  Then
\begin{equation}
 \frac{\sqrt\kappa+\kappa}{8A}
 \le Q^*(\kappa,A)
 \le\frac{1+2\kappa}{A\sqrt{1+\kappa^{-1}}}
 \le\frac{3(\sqrt\kappa+\kappa)}A.
 \label{eq:scalar-rms-amplitude-law}
\end{equation}
For rational $\kappa,A$, the upper construction has expected auxiliary
bit work
\[
 O\!\left(B^4\left[1+\frac{\sqrt\kappa+\kappa}{A}\right]\right),
\]
where $B\ge16$ includes their encoding lengths and source-address
lengths.  Its expected source count has the explicit middle bound in
\eqref{eq:scalar-rms-amplitude-law}.
\end{theorem}
\paragraph{Proof idea.}
A known thinning gate rescales the geometric-majority identity
to the requested amplitude. Curvature gives the lower bound when the
normalization is steep; coupling the two deterministic source biases
handles the remaining parameter range.

\begin{proof}
Put $\theta=\kappa^{-1}$ and
$q=1/[A\sqrt{1+\theta}]\le1$. Lemma~\ref{lem:rmsgeommajority} gives
\[
 \frac{z\sqrt{1+\theta}}{\sqrt{\theta+z^2}}
 =\mathbb E M_K(z),\qquad
 \Pr(K=k)=\frac\theta{1+\theta}
               \left(\frac1{1+\theta}\right)^k,
 \quad k\ge0.
\]
First open a known gate of probability $q$.  On its closed branch
return a fair sign.  On its open branch draw $K$ and return the
majority of $2K+1$ fresh source signs.  This yields exactly
$f_{\kappa,A}$.  Since $\mathbb EK=\kappa$, its complete-majority
source cost is $q(1+2\kappa)$.  If $\kappa\le1$, use
$q\le\sqrt\kappa/A$ and $1+2\kappa\le3$.  If $\kappa\ge1$, use
$q\le1/A$ and $1+2\kappa\le3\kappa$.  These prove the upper bounds.

For $\kappa\ge1$, direct differentiation gives
\[
 |f_{\kappa,A}''(t)|
 =\frac{3\kappa^{-1}|t|}
        {A(\kappa^{-1}+t^2)^{5/2}}.
\]
Take $t=1/(2\sqrt\kappa)$.  Then
\[
 |f_{\kappa,A}''(t)|=\frac{48\kappa}{25\sqrt5 A},\qquad
 1-t^2\ge\frac34.
\]
Lemma~\ref{lem:new-curvature-transfer} gives
\[
 Q^*(\kappa,A)\ge\frac{18\kappa}{25\sqrt5 A}
                 >\frac\kappa{4A}
                 \ge\frac{\sqrt\kappa+\kappa}{8A}.
\]
The strict comparison follows by squaring $72>25\sqrt5$.

For $0<\kappa\le1$, run the factory at source means $+1$ and $-1$.
Every source request then returns the same deterministic sign, so the
factory reads one input coordinate in $\{-1,1\}$.
Lemma~\ref{lem:attentioncoupling} bounds the probability of any
source request below by the total variation distance between the
two required output laws,
\[
 f_{\kappa,A}(1)=\frac{\sqrt\kappa}{A\sqrt{1+\kappa}}
                \ge\frac{\sqrt\kappa}{\sqrt2 A}.
\]
Every uniform expected-cost bound is at least this probability,
which is more than the stated lower bound.

The outer known gate uses certified square-root intervals and a lazy
uniform comparison, with expected $O(B^4)$ auxiliary work.  Draw $K$
only after that gate opens.  The simplest exact implementation flips the known rational coin
with success probability $1/(1+\kappa)$ until its first success
and lets $K$ count the failures.  Each such coin has expected
$O(B^2)$ bit work by integer rejection.  Its geometric tail gives,
for each fixed integer $r$, a bound
\[
 \mathbb E[(1+K)(1+\log(K+2))^r]
 \le C_r(1+\kappa)(1+\log(\kappa+2))^r.
\]
This follows by dividing the geometric distribution into intervals
of length $\lceil1+\kappa\rceil$ and summing its exponential tail.
Counters, source addresses, and majority evaluation therefore have
conditional expected auxiliary cost $O((1+\kappa)B^4)$.  The initial
gate multiplies this cost by $q$.  The preceding inequalities finish
the bit bound.  Known comparisons stop almost surely, $K$ is finite
almost surely, and the finite majority terminates.
\end{proof}

For the earlier scalar floor interface, let $\kappa=R^2/\epsilon\ge1$
and return RMS at radius $2R/\sqrt\epsilon=2\sqrt\kappa$.
This is $A=2\sqrt\kappa$ in the theorem, so its optimal uniform
source cost has order $\sqrt\kappa$.  At the fixed radius $A=4$, the
cost instead has order $\kappa$.  Both statements are consistent;
they request different output amplitudes.
